% !TeX root = ../Cosmology.tex

%% --------------------------------------------------------
\chapter{Програмні пакети Python для космологічних досліджень}
\label{sec:python_packages}
%% --------------------------------------------------------

Сучасна теоретична та спостережна космологія є високотехнологічною дисципліною, де перевірка теоретичних моделей та аналіз даних експериментів (таких як \emph{Planck}, \emph{Euclid} чи \emph{JWST}) спираються на потужний інструментарій чисельного моделювання. Найбільшого поширення в академічному середовищі набула екосистема мови Python, яка об'єднує спеціалізовані космологічні пакети та загальнонаукові бібліотеки для аналізу даних і машинного навчання.

Нижче наведено та класифіковано основні програмні пакети, які є стандартом у сучасних космологічних дослідженнях:

\begin{enumerate}[leftmargin=*]
    \item \textbf{Пакети для лінійного аналізу збурень та еволюції СМВ:}
    \begin{itemize}
        \item \href{http://class-code.net/}{\texttt{CLASS}} (\emph{Cosmic Linear Anisotropy Solving System}) --- один із двох найпопулярніших кодів для інтегрування лінійних збурень у ранньому Всесвіті. Написаний на мові C, але має повну та зручну обгортку \texttt{classy} для Python. Пакет дозволяє з високою точністю розраховувати спектри потужності температурних анізотропій та поляризації реліктового випромінювання, а також спектри потужності матерії $P(k)$.
        \item \href{https://camb.readthedocs.io/}{\texttt{CAMB}} (\emph{Code for Anisotropy in the Microwave Background}) --- класичний інструмент для розрахунку анізотропії CMB. Його обчислювальне ядро реалізоване на Fortran, проте Python-версія повністю інтегрована в сучасні конвеєри обробки даних і дозволяє гнучко модифікувати рівняння стану темної енергії або параметри первинного спектра збурень.
    \end{itemize}

    \item \textbf{Статистичний аналіз та оцінка параметрів (MCMC sampling):}
    \begin{itemize}
        \item \href{https://cobaya.readthedocs.io/}{\texttt{Cobaya}} --- сучасний фреймворк для байєсівського аналізу космологічних даних, розроблений для роботи з архівами \emph{Planck} та майбутніх місій. \texttt{Cobaya} дозволяє зв'язувати розрахунки з \texttt{CAMB} або \texttt{CLASS} та запускати алгоритми Монте-Карло за схемами марковських ланцюгів (MCMC) для пошуку обмежень на космологічні параметри ($\Omega_m$, $H_0$, $n_s$ тощо).
        \item \href{https://github.com/brinckmann/montepython_public}{\texttt{MontePython}} --- популярний MCMC-семплер, який традиційно використовується у зв'язці з кодом \texttt{CLASS} для порівняння теоретичних моделей із космологічними спостереженнями та побудови контурів довірчих інтервалів.
    \end{itemize}

    \item \textbf{Загальні астрофізичні фреймворки та обробка даних:}
    \begin{itemize}
        \item \href{https://www.astropy.org/}{\texttt{Astropy}} --- фундаментальна бібліотека для астрономічних досліджень. Модуль \texttt{astropy.cosmology} містить готові класи для швидкого розрахунку різних типів відстаней (супутньої, кутової, світності) та еволюції густини енергії для стандартних моделей ($\Lambda\text{CDM}$, $w\text{CDM}$).
        \item \href{https://healpy.readthedocs.io/}{\texttt{healpy}} --- Python-обгортка для бібліотеки HEALPix (\emph{Hierarchical Equal Area isoLatitude Pixelization}). Це критично важливий інструмент для роботи з картами реліктового випромінювання та великомасштабної структури, розбитими на сфері. Пакет дозволяє виконувати сферичне гармонічне перетворення для отримання мультипольних коефіцієнтів $C_\ell$.
    \end{itemize}
\end{enumerate}

Для чисельного розв'язання нестандартних динамічних систем (наприклад, для точного інтегрування рівняння Кляйна-Ґордона-Фока зі складними видами потенціалів інфлатона на етапі розігріву) зазвичай використовуються базові наукові пакети \href{https://numpy.org/}{\texttt{NumPy}} та \href{https://scipy.org/}{\texttt{SciPy}} (зокрема, модулі для розв'язання жорстких систем звичайних диференціальних рівнянь \texttt{scipy.integrate}). Візуалізація отриманих спектрів та еволюційних треків традиційно виконується за допомогою бібліотеки \href{https://matplotlib.org/}{\texttt{Matplotlib}}.


%% --------------------------------------------------------
\section{Спостережні дані космічних місій на GitHub}
%% --------------------------------------------------------

Сучасні космологічні дослідження неможливі без доступу до реальних
спостережних даних. Нижче наведено ключові репозиторії, де зберігаються
дані основних космічних місій у форматі, готовому для використання з
фреймворками типу \texttt{Cobaya} та \texttt{MontePython}.

\subsection*{Planck (2018 release, PR3/PR4)}

Офіційні дані місії \emph{Planck} є золотим стандартом за посиланням на аналіз CMB.
Репозиторії містять як спектри потужності, так і \texttt{likelihood} для
MCMC-аналізу.

\begin{itemize}
    \item \href{https://github.com/heatherprince/planck-lite-py}{\texttt{planck-lite-py}} ---
          Python-реалізація спрощеного \texttt{plik-lite} likelihood для
          Planck 2015/2018. Дозволяє швидко обчислювати $\chi^2$ для
          TT, TE, EE спектрів. \cite{prince_dunkley_2019}

    \item \href{https://github.com/heatherprince/planck-low-py}{\texttt{planck-low-py}} ---
          Стиснутий низькомультипольний ($\ell<30$) likelihood для
          температури та $E$-мод поляризації. Оптимальний для спільних
          аналізів із іншими експериментами. \cite{prince_low_2021}

    \item \href{https://github.com/benabed/clipy}{\texttt{clipy}} ---
          Чиста Python-реалізація офіційного Planck \texttt{clik} з
          підтримкою \texttt{jax}. Сумісна з \texttt{plik} (високі $\ell$),
          \texttt{simall} (низькі $\ell$ поляризації) та
          \texttt{commander} (низькі $\ell$ температура). \cite{benabed_2025}

    \item \href{https://github.com/carronj/planck_PR4_lensing}{\texttt{planck\_PR4\_lensing}} ---
          Дані гравітаційного лінзування CMB за результатами
          Planck PR4 (NPIPE). Включає карти збіжності лінзування
          (TT, polarization-only, MV) та MCMC-ланцюжки. \cite{carron_2022}

    \item \href{https://github.com/CobayaSampler/planck_lensing_external}{\texttt{planck\_lensing\_external}} ---
          Приклад зовнішнього likelihood для лінзування Planck 2018,
          сумісний з \texttt{Cobaya}.
\end{itemize}

\subsection*{CMB-S4, SPT-3G, Advanced SO}

Для майбутніх та діючих експериментів існують \emph{mock} дані, які
використовуються для прогнозування точності майбутніх вимірювань.

\begin{itemize}
    \item \href{https://github.com/sriniraghunathan/CMB_BAO_SNe_likelihoods}{\texttt{CMB\_BAO\_SNe\_likelihoods}} ---
          \texttt{Cobaya}-сумісні \emph{mock} likelihood для спільного
          аналізу CMB, BAO та SNe. Включає:
          \begin{itemize}
              \item Mock CMB дані для \textbf{SPT-3G}, \textbf{Advanced SO-Baseline},
                    \textbf{ASO-Goal} та \textbf{CMB-S4}.
              \item Mock SNe дані для \textbf{DES} та \textbf{LSST}.
              \item Mock BAO дані для \textbf{DESI-DR2} та \textbf{DESI-DR3}.
          \end{itemize}
          Цей репозиторій є чудовим інструментом для \emph{forecasting} ---
          оцінки того, наскільки точно майбутні місії зможуть виміряти
          космологічні параметри. \cite{raghunathan_2025}
\end{itemize}

\subsection*{BAO та SNe дані}

\begin{itemize}
    \item \textbf{DESI} (Dark Energy Spectroscopic Instrument) ---
          дані баріонних акустичних осциляцій доступні в репозиторіях,
          інтегрованих у \texttt{Cobaya} (DR2, DR3).

    \item \textbf{Pantheon+} --- найбільша збірка кривих блиску наднових
          типу Ia, доступна для аналізу в \texttt{Cobaya}.
\end{itemize}

\subsection*{Вимірювання $H_0$ (Hubble tension)}

Для аналізу напруженості Хаббла \texttt{Cobaya} надає вбудовані
гауссові \texttt{likelihood} для локальних вимірювань $H_0$:
\begin{itemize}
    \item \texttt{H0.riess2018a}, \texttt{H0.riess2018b} --- Riess et al. 2018.
    \item \texttt{H0.riess201903} --- Riess et al. 2019.
    \item \texttt{H0.riess2020} --- Riess et al. 2020.
    \item \texttt{H0.freedman2020} --- Freedman et al. 2020.
    \item \texttt{H0.riess2020Mb} --- обмеження на абсолютну зоряну
          величину $M_b$ (для сумісного аналізу з SNe).
\end{itemize}
Документацію див. у \href{https://cobaya.readthedocs.io/en/latest/likelihood_H0.html}{Cobaya docs}.

\subsection*{Як отримати дані безпосередньо з коду}

Більшість перелічених репозиторіїв можна клонувати безпосередньо:
\begin{verbatim}
git clone https://github.com/heatherprince/planck-lite-py.git
git clone https://github.com/sriniraghunathan/CMB_BAO_SNe_likelihoods.git
\end{verbatim}

Для офіційних даних Planck (спектри, карти) використовуйте
\href{https://pla.esac.esa.int/pla}{Planck Legacy Archive (PLA)}.
Дані PR4 (NPIPE) також доступні через
\href{https://data.cmb-s4.org/planck_pr4.html}{CMB-S4 Data Portal}.